\section{Experiments}
\label{sec:experiments}
We evaluate the proposed framework by posing three questions.
\textbf{Q1}, Does the variable-cardinality deferral improve the trade-off between segmentation quality and expert cost?
\textbf{Q2}, Does the proposed multi-frame selection model identify better correction frames than existing frame-selection methods under a fixed correction budget?
\textbf{Q3}, Does the cardinality model allocate expert effort to videos that benefit most from additional corrections?

\subsection{Experimental setup}
\label{sec:experimental_setup}

\paragraph{Datasets and prompting protocol.}
We evaluate on VOST \cite{tokmakov2023breaking}, CAMotion
\cite{yao2026camotion}, SUN-SEG \cite{ji2022vps}, and VTUS
\cite{yang2024vivim}, using the official train/test splits.
We sample ten uniformly spaced frames per video for prompting, with an initial ground-truth mask provided in the first frame to identify the object of interest. This sampling protocol enables tractable exhaustive oracle comparisons, while our framework operates without combinatorial search. We further evaluate larger candidate pools in \ref{sec:ablation}.
SAM2 propagates over the complete video, and training rewards
and video-level mIoU are computed over all frames.
Unless otherwise stated, corrective prompts are ground-truth
masks, and we allow up to $K=4$ 
% \lokesha{make this 3, if llm flaged 10 frame limitation we need to handle it}
%\lokesha{later update mask to point prompt} 
additional corrections,
enabling the method to address multiple propagation failures
while evaluating variable-cardinality decisions within a
bounded correction budget.
Additional details and results for other values of $K$
are provided in the appendix.

\paragraph{Implementation details.}
We use a frozen SAM2.1 Hiera-Tiny model \cite{ravi2025sam} as the
promptable video segmentation model. For each sampled frame, we use
embeddings produced by the SAM2 memory encoder, which jointly encode
frame appearance and the corresponding predicted mask.
These representations are processed by independently trained
spatial-temporal encoders for the multi-frame selection and
cardinality models, followed by MLP decision heads.
The frame selector additionally uses candidate-to-context attention
and learned cardinality embeddings.

The frame selector is trained for 400 epochs with $B=4$ sampled
alternatives, after which it is frozen and the cardinality model
is trained for 200 epochs. Both models use AdamW (learning rate
$10^{-5}$, weight decay $10^{-4}$, and batch size 16) and are trained
on an NVIDIA GeForce RTX 5090 GPU.
Unless otherwise stated, we use $\alpha=0.01$ and train separate
cardinality models for different expert-cost weights to evaluate
the quality--cost trade-off.
All learned methods are evaluated over three random seeds.
Reported $\pm$ values denote the standard deviation across these runs.
Additional implementation details are provided in the appendix.
%\lokesha{here we need to update architecture details and remove about B and ek vector, do wonder whter we need std here, or either confidence interval}

\begin{table}[t]
\caption{Fixed-cardinality frame selection at \(k=3\). Rounds column shows $1$ when all three frames are selected in a single round, $3$ when one frame is selected per round.}
\label{tab:fixed_cardinality}

\centering
\small
\setlength{\tabcolsep}{2.4pt}
\renewcommand{\arraystretch}{1.0}

\resizebox{\columnwidth}{!}{
\begin{tabular}{@{}llcccc@{}}
\toprule
\textbf{Method} &
\textbf{Rounds} &
\textbf{VOST} &
\textbf{CAMotion} &
\textbf{SUN-SEG} &
\textbf{VTUS}
\\
\midrule
Random frames      & 1      & 0.613 & 0.831 & 0.787 & 0.727 \\
Min-IoU frames   & 1      & 0.647 & 0.843 & 0.791 & 0.748 \\
HVSA \cite{hvsa}        & 1      & 0.663 & 0.844 & 0.828 & 0.781 \\
IVOS-W \cite{IVOSW}      & 3 & 0.655 & 0.845 & 0.818 & 0.769 \\
EVA-VOS \cite{delatolas2024learning}     & 3 & 0.667 & 0.848 & 0.825 & 0.786 \\
Top-$k$ \cite{montreuil2025ask}     & 1      & 0.670 & 0.847 & 0.829 & 0.772 \\
\textbf{Ours}      & \textbf{1}
                   & \textbf{0.688}
                   & \textbf{0.852}
                   & \textbf{0.862}
                   & \textbf{0.799} \\
\midrule
\textit{Oracle}    & \textit{1}
                   & \textit{0.712}
                   & \textit{0.859}
                   & \textit{0.875}
                   & \textit{0.828} \\
\bottomrule
\end{tabular}
}


\end{table}


\paragraph{Evaluation metrics.}
We report video-level mean intersection over union (mIoU), the mean number of requested corrections ($\bar{k}$), and the average system cost, defined as $1-\mathrm{mIoU}+\alpha\bar{k}$, where $\alpha$ is a fixed
penalty assigned to each requested correction.
%\lokesha{say statisttical details , while main evaluation metrix is iou JandF are in supplement}



\paragraph{Evaluation protocols and comparison methods.}


\noindent\textbf{Q1: How effectively does the proposed framework
balance segmentation quality and expert correction cost?}
We evaluate variable-cardinality deferral under a maximum correction
budget of $K=4$, comparing against a no-correction reference,
a random policy, and four baselines.
Lazy-XMem~\cite{vujasinovic2024strike} triggers corrections
using per-frame uncertainty, while
L2RP~\cite{rasanjalee2026understanding} supports only
$k\in\{0,1\}$. We additionally adapt two cost-aware learning-to-defer (L2D)
approaches originally proposed for multiple-expert classification
to our video frame-selection setting, treating candidate correction
frames as expert actions. LECODU~\cite{zhang2024learning}
predicts the number of corrections to request and selects the
corresponding frames at random, while Adaptive Top-$k(x)$~
\cite{montreuil2025ask} predicts the correction cardinality
and selects frames according to independently estimated
correction utilities. We also include an oracle upper bound
that jointly selects the cardinality and correction frames
to minimise system cost.

We further examine the quality--cost trade-off across different
expert-cost weights $\alpha$. To analyse deferral decisions,
we evaluate mIoU across non-deferral coverage levels by varying
the deferral threshold. Finally, we analyse the adaptive allocation of corrections
by grouping videos according to their initial segmentation
mIoU and examining the mean selected cardinality and
resulting mIoU improvement across different quality levels.

\noindent\textbf{Q2: How effective is the proposed multi-frame
selection model under a fixed correction budget?}
To evaluate frame selection independently of cardinality decisions,
we fix the number of corrections per video and compare against
random selection, a minimum-IoU heuristic that selects frames
with the lowest IoU between the initial propagation and ground truth,
a fixed-cardinality variant of Top-$k(x)$~\cite{montreuil2025ask}
that selects the $k$ frames with the highest independently estimated
correction utilities, and adaptations of HVSA~\cite{hvsa},
EVA-VOS~\cite{delatolas2024learning}, and IVOS-W~\cite{IVOSW}.
HVSA and Top-$k(x)$ select correction frames in a batch, whereas
EVA-VOS and IVOS-W select frames sequentially, re-propagating
after each correction. Our method selects all correction frames
in a single batch without intermediate re-propagation.
We also report an oracle upper bound that selects the frame
subset yielding the highest mIoU for each fixed $k$.

We additionally present ablation studies, robustness evaluations,
and qualitative results. Further details of the baseline adaptations,
evaluation protocols, and additional analyses are provided in the appendix.



